%! TEX root = main.tex

\section{Modified Pressure Flux}
\label{sc:33}

The standard SWE system \eqref{eq:swe_conservative} assumes a hydrostatic pressure law of the form \( p(h) = \tfrac{1}{2}gh^2 \), which leads to the quadratic pressure term in the second component of the flux vector \eqref{eq:swe_vectors}. This section considers a modification in which the pressure law is replaced by a logarithmic form, giving rise to an alternative conservative system whose mathematical structure and wave dynamics differ in a physically meaningful way from the standard model.

\subsection{Analytical Formulation}
\label{ssc:331}

The modified pressure law is defined as

\begin{equation}
p(h) = gh \log\left(\frac{h}{h_{\mathrm{ref}}}\right)
\label{eq:mod_pressure}
\end{equation}

where \( h_{\mathrm{ref}} > 0 \) is a fixed reference height. Substituting \eqref{eq:mod_pressure} into the momentum equation in place of the standard hydrostatic term yields the modified flux vector

\begin{equation}
\underline{F}(\underline{U}) =
\begin{bmatrix}
q \\
\dfrac{q^2}{h} + gh\log\left(\dfrac{h}{h_{\mathrm{ref}}}\right)
\end{bmatrix}
\label{eq:mod_flux}
\end{equation}

while the mass conservation equation and the state vector \( \underline{U} = [h, q]^T \) remain unchanged. The system retains the conservative form \eqref{eq:swe_conservative} with the flux \eqref{eq:mod_flux}, and the same FVM update formula \eqref{eq:fvm_update} applies without modification: only the evaluation of the flux vector at each interface changes. The mathematical properties of the modified system differ from those of the standard SWE because the Jacobian \( \mathbf{J} = \partial \underline{F} / \partial \underline{U} \) now takes the form

\begin{equation}
\mathbf{J}(\underline{U}) =
\begin{bmatrix}
0 & 1 \\
-\dfrac{q^2}{h^2} + g\log\left(\dfrac{h}{h_{\mathrm{ref}}}\right) + g & \dfrac{2q}{h}
\end{bmatrix}
\label{eq:mod_jacobian}
\end{equation}

The eigenvalues of \( \mathbf{J} \), which govern the characteristic speeds of the system, are

\begin{equation}
\lambda_{1,2} = u \mp c_{\mathrm{mod}}, \qquad c_{\mathrm{mod}} = \sqrt{g\left[\log\left(\frac{h}{h_{\mathrm{ref}}}\right) + 1\right]}
\label{eq:mod_charspeed}
\end{equation}

where \( c_{\mathrm{mod}} \) plays the role of the wave celerity for the modified system. Comparing with the standard celerity \( c = \sqrt{gh} \), the modified celerity \( c_{\mathrm{mod}} \) depends logarithmically rather than algebraically on \( h \), which has several physical consequences. For \( h > h_{\mathrm{ref}}  e^{-1} \) the argument of the square root in \eqref{eq:mod_charspeed} is positive and the system remains hyperbolic; below this threshold the characteristic speeds become complex, signaling a loss of hyperbolicity that renders the initial-value problem ill-posed.

\paragraph{Critical remark on the dam-break test.}
In the dam-break initial data (case~a, section~\ref{ssc:321}), 
the right state has $h_R = 0.1\,\mathrm{m}$. Since
%
\begin{equation}
    h_R = 0.1 < h_{\mathrm{ref}}\,e^{-1} = e^{-1} \approx 0.368\,\mathrm{m}
    \label{eq:hyp_threshold}
\end{equation}
%
the modified system \eqref{eq:mod_flux} is \textbf{ill-posed} 
($c_{\mathrm{mod}}^2 < 0$) in the right half of the domain at 
$t = 0$. The Rusanov flux remains formally applicable in this 
region but the numerical diffusion  artificially stabilises what is in fact an ill-posed problem, and the computed solution there lacks physical justification. The results of figure ~\ref{fig:mod_dam_break} for the modified-pressure dam-break must therefore be interpreted with  caution: they are numerically stable but not physically meaningful in the region $h < h_{\mathrm{ref}}\,e^{-1}$.

%In the hyperbolic regime, the slower growth of \( c_{\mathrm{mod}} \) with \( h \) compared to \( \sqrt{gh} \) implies that the modified system propagates waves more slowly at large depths and produces weaker shock strengths for the same initial height discontinuity. The rarefaction fans are correspondingly broader, and the overall flow is less compressive than in the standard shallow-water setting. These differences in characteristic structure are directly observable in the numerical solutions discussed below.


\subsection{Numerical Results and Discussion}
\label{ssc:332}

The modified flux \eqref{eq:mod_flux} is tested with \( h_{\mathrm{ref}} = 1 \) on the dam-break and stationary-shock initial data, using the first-order Godunov scheme for both the standard and modified systems to isolate the effect of the pressure law from any scheme-dependent dissipation or dispersion.

\paragraph{Dam break with modified pressure}
Figure \ref{fig:mod_dam_break} compares the two Godunov solutions at time \( T = \SI{1}{\second} \) for the dam-break initial data (\( h_L = 1 \), \( h_R = 0.1 \)). The most immediate effect of the modified pressure is a marked reduction in the rightward shock speed: at \( T = 1 \) the shock front computed with the standard flux has already reached and reflected from the right wall, whereas the modified-pressure shock is still propagating through the interior of the domain, having traveled a substantially shorter distance. This behavior is in direct agreement with the eigenvalue analysis of section \ref{ssc:331}: for \( h \leq 1 \) and \( h_{\mathrm{ref}} = 1 \) one has \( \log(h/h_{\mathrm{ref}}) \leq 0 \), so that \( c_{\mathrm{mod}} \leq \sqrt{g} \) is strictly smaller than the corresponding standard celerity \( c = \sqrt{gh} \) evaluated at the same depth. The weaker pressure gradient of the logarithmic law provides less driving force per unit height, resulting in a slower propagation of the shock front and a less pronounced post-shock height plateau. The rarefaction fan on the left is also affected: the modified solution shows a smoother and more extended depletion profile rather than the flat constant-state region produced by the standard flux, reflecting the altered characteristic fan structure. The discharge panel confirms both observations: the peak value of \( q \) is significantly lower in the modified case, and the profile shape is qualitatively different, with a more gradual spatial variation across the entire domain.

\begin{figure}[!htb]
    \centering
    \includegraphics[width=0.90\textwidth]{00_Images/HW3/3.4_ModifiedDamBreak.png}
    \caption{Dam-break Riemann problem at \(T = \SI{1}{\second}\): comparison between the standard Godunov scheme (solid blue) and the modified-pressure Godunov scheme with \(h_{\mathrm{ref}} = 1\) (dashed orange). Left panel: water height \(h(x,T)\). The modified shock front has not yet reached the right wall, travelling at a lower speed consistent with the reduced wave celerity \(c_{\mathrm{mod}}\) predicted by \eqref{eq:mod_charspeed}; the post-shock plateau is lower and the rarefaction profile is more gradual. Right panel: discharge \(q(x,T)\); the peak discharge is substantially reduced in the modified case, reflecting the weaker momentum transfer induced by the logarithmic pressure term. Physical parameters: \(L = \SI{1}{\meter}\), \(g = \SI{9.81}{\meter\per\second\squared}\), \(h_L = 1\), \(h_R = 0.1\), \(h_{\mathrm{ref}} = 1\)}
    \label{fig:mod_dam_break}
\end{figure}

\paragraph{Stationary shock with modified pressure}
The most striking consequence of the pressure modification is revealed by the stationary-shock test, shown in figure \ref{fig:mod_stationary_shock}. In the standard system, the initial data \( (h_L, h_R, q) \) are constructed precisely so that the Rankine-Hugoniot jump condition yields a shock speed \( \sigma = 0 \): the net flux imbalance across the discontinuity vanishes, \( [\underline{F}_{\mathrm{std}}] = \underline{0} \), and the shock remains stationary at the center of the domain for all time. Replacing the standard flux with the modified flux \eqref{eq:mod_flux} alters the second component of the flux vector at every height level, so that the same left and right states now produce a non-zero flux jump \( [\underline{F}_{\mathrm{mod}}] \neq \underline{0} \). 

The Rankine-Hugoniot condition for the modified system reads

\begin{equation}
    \sigma[\underline{U}] = [\underline{F}_{\mathrm{mod}}(\underline{U})]
    \label{eq:rh_mod}
\end{equation}

The shock speed is obtained by combining the two scalar components 
of \eqref{eq:rh_mod}. Eliminating $[q]$ between them gives 
$\sigma^2 = [\underline{F}_{\cdot,2}]/[h]$, so:

\begin{equation}
    \sigma_{\mathrm{std}}^2 = \frac{\tfrac{1}{2}g[h^2]}{[h]} 
    = \tfrac{1}{2}g(h_L+h_R)
    \label{eq:sigma_std_explicit}
\end{equation}

\begin{equation}
    \sigma_{\mathrm{mod}}^2 
    = \frac{g\!\left[h\log\!\left(\dfrac{h}{h_{\mathrm{ref}}}
      \right)\right]}{[h]}
    = \frac{9.81\times\bigl[0.5\log(0.5) - \log(1)\bigr]}{-0.5}
    = 6.80\;\mathrm{m^2/s^2}
    \label{eq:sigma_mod_explicit}
\end{equation}

\begin{equation}
    \sigma_{\mathrm{std}} \approx 2.71\;\mathrm{m/s} 
    > \sigma_{\mathrm{mod}} \approx 2.61\;\mathrm{m/s}
    \label{eq:sigma_comparison}
\end{equation}

The logarithmic pressure law produces a \textbf{weaker shock speed} than the standard SWE for the same initial height jump, consistent with the eigenvalue comparison $c_{\mathrm{mod}} < c$ established in section~\ref{ssc:331}. The Rankine-Hugoniot condition \eqref{eq:rh_mod} is therefore satisfied only for a strictly positive shock speed \( \sigma > 0 \), meaning that the hydraulic jump is no longer stationary but drifts rightward at a finite velocity. This is precisely what figure~\ref{fig:mod_stationary_shock} shows: at \( T = \SI{1}{\second} \) the modified pressure shock has migrated visibly to the right of the standard one, while both solutions retain a sharp, monotone profile characteristic of the first-order Godunov scheme. The pre-shock heights on the left plateau also differ slightly between the two cases, reflecting the different intermediate states produced by the respective Riemann solvers. The discharge panel corroborates the picture: the modified solution exhibits a higher and spatially shifted peak, consistent with a shock that is still actively propagating and carrying momentum to the right rather than being in a force-balanced stationary configuration.

\begin{comment}
The shock speed is obtained by combining the two Rankine-Hugoniot 
equations \eqref{eq:rh_mod}. Eliminating $[q]$ between them gives 
$\sigma^2 = [\underline{F}_{\cdot,2}]/[h]$, so:
%
\begin{equation}
    \sigma_{\mathrm{std}}^2 = \frac{\tfrac{1}{2}g[h^2]}{[h]} 
    = \tfrac{1}{2}g(h_L+h_R)
    \label{eq:sigma_std_explicit}
\end{equation}
%
\begin{equation}
    \sigma_{\mathrm{mod}}^2 
    = \frac{g\!\left[h\log\!\left(\dfrac{h}{h_{\mathrm{ref}}}\right)\right]}{[h]}
    = \frac{9.81\times\bigl[0.5\log(0.5) - \log(1)\bigr]}{-0.5}
    = 6.80\;\mathrm{m^2/s^2}
    \label{eq:sigma_mod_explicit}
\end{equation}
%
\begin{equation}
    \sigma_{\mathrm{std}} \approx 2.71\;\mathrm{m/s} 
    > \sigma_{\mathrm{mod}} \approx 2.61\;\mathrm{m/s}
    \label{eq:sigma_comparison}
\end{equation}


The shock speed $\sigma_{\mathrm{mod}}$ can be estimated by evaluating 
the second component of the jump condition \eqref{eq:rh_mod} explicitly. 
With $q_L = q_R = 0$, the mass equation gives $\sigma[h] = 0$, which 
forces the shock to carry a net mass flux; the momentum equation yields

\begin{equation}
    \sigma_{\mathrm{mod}} 
    = \frac{[\,\underline{F}_{\mathrm{mod},2}\,]}{[\,q\,]}
    = \frac{g\!\left[h\log\!\left(\dfrac{h}{h_{\mathrm{ref}}}\right)
      \right]}{[q]}
    \label{eq:sigma_mod_explicit}
\end{equation}

For the standard system the analogous expression involves 
$\tfrac{1}{2}g[h^2]$ in the numerator. Evaluating both pressure 
jumps with $h_L = 1$, $h_R = 0.5$, $h_{\mathrm{ref}} = 1$:

\begin{align}
    \text{Standard:} \quad 
    \tfrac{1}{2}g[h^2] 
        &= \tfrac{1}{2}\times 9.81\times(0.25 - 1) 
         = -3.68\;\mathrm{m^2/s^2}
    \label{eq:pressure_jump_std}\\[4pt]
    \text{Modified:} \quad 
    g\!\left[h\log\!\tfrac{h}{h_{\mathrm{ref}}}\right] 
        &= 9.81\times\bigl[0.5\log(0.5) - 1\cdot\log(1)\bigr]
         = 9.81\times(-0.347)
         \approx -3.40\;\mathrm{m^2/s^2}
    \label{eq:pressure_jump_mod}
\end{align}


The magnitude of the modified pressure jump ($3.40\,\mathrm{m^2/s^2}$) is smaller than the standard one ($3.68\,\mathrm{m^2/s^2}$). Since the denominator $[q]$ is driven by the same initial conditions in both  cases, it follows that $|\sigma_{\mathrm{mod}}| < |\sigma_{\mathrm{std}}|$: the logarithmic pressure law produces a \textbf{weaker pressure  imbalance} across the discontinuity, driving the shock at a lower speed.  This is consistent with the eigenvalue comparison $c_{\mathrm{mod}} < c$  established in section~\ref{ssc:331} and directly visible in  figure~\ref{fig:mod_stationary_shock}.

is therefore satisfied only for a strictly positive shock speed \( \sigma > 0 \), meaning that the hydraulic jump is no longer stationary but drifts rightward at a finite velocity. This is precisely what figure \ref{fig:mod_stationary_shock} shows: at \( T = \SI{1}{\second} \) the modified-pressure shock has migrated visibly to the right of the standard one, while both solutions retain a sharp, monotone profile characteristic of the first-order Godunov scheme. The pre-shock heights on the left plateau also differ slightly between the two cases, reflecting the different intermediate states produced by the respective Riemann solvers. The discharge panel corroborates the picture: the modified solution exhibits a higher and spatially shifted peak, consistent with a shock that is still actively propagating and carrying momentum to the right rather than being in a force-balanced stationary configuration.
\end{comment}

\begin{figure}[!htb]
    \centering
    \includegraphics[width=0.90\textwidth]{00_Images/HW3/3.5_ModifiedShock.png}
    \caption{Stationary-shock Riemann problem at \(T = \SI{1}{\second}\): comparison between the standard Godunov scheme (solid blue) and the modified-pressure Godunov scheme with \(h_{\mathrm{ref}} = 1\) (dashed orange). Left panel: water height \(h(x,T)\). The initial data were tuned so that the standard Rankine-Hugoniot condition yields \(\sigma = 0\); the modified pressure alters the flux balance according to \eqref{eq:rh_mod}, producing a positive shock speed that causes the hydraulic jump to drift rightward, as clearly visible in the shifted front position. Right panel: discharge \(q(x,T)\); the higher peak and rightward displacement of the modified solution confirm the active momentum transport induced by the non-zero shock speed. Physical parameters: \(L = \SI{1}{\meter}\), \(g = \SI{9.81}{\meter\per\second\squared}\), \(h_L = 1\), \(h_R = 0.5\), \(h_{\mathrm{ref}} = 1\)}
    \label{fig:mod_stationary_shock}
\end{figure}